\documentclass[a4paper,11pt,oneside]{memoir}

\usepackage[table,xcdraw]{xcolor}

% Castellano
\usepackage[spanish,es-tabla]{babel}
\selectlanguage{spanish}
\usepackage[utf8]{inputenc}
\usepackage{placeins}

\RequirePackage{booktabs}
\RequirePackage[table]{xcolor}
\RequirePackage{xtab}
\RequirePackage{multirow}

% Links
\usepackage[colorlinks]{hyperref}
\hypersetup{
	allcolors = {red}
}

% Ecuaciones
\usepackage{amsmath}

% Rutas de fichero / paquete
\newcommand{\ruta}[1]{{\sffamily #1}}

% Párrafos
\nonzeroparskip


% Imagenes
\usepackage{graphicx}
\newcommand{\imagen}[2]{
	\begin{figure}[!h]
		\centering
		\includegraphics[width=0.9\textwidth]{#1}
		\caption{#2}\label{fig:#1}
	\end{figure}
	\FloatBarrier
}

\newcommand{\imagenflotante}[2]{
	\begin{figure}%[!h]
		\centering
		\includegraphics[width=0.9\textwidth]{#1}
		\caption{#2}\label{fig:#1}
	\end{figure}
}



% El comando \figura nos permite insertar figuras comodamente, y utilizando
% siempre el mismo formato. Los parametros son:
% 1 -> Porcentaje del ancho de página que ocupará la figura (de 0 a 1)
% 2 --> Fichero de la imagen
% 3 --> Texto a pie de imagen
% 4 --> Etiqueta (label) para referencias
% 5 --> Opciones que queramos pasarle al \includegraphics
% 6 --> Opciones de posicionamiento a pasarle a \begin{figure}
\newcommand{\figuraConPosicion}[6]{%
  \setlength{\anchoFloat}{#1\textwidth}%
  \addtolength{\anchoFloat}{-4\fboxsep}%
  \setlength{\anchoFigura}{\anchoFloat}%
  \begin{figure}[#6]
    \begin{center}%
      \Ovalbox{%
        \begin{minipage}{\anchoFloat}%
          \begin{center}%
            \includegraphics[width=\anchoFigura,#5]{#2}%
            \caption{#3}%
            \label{#4}%
          \end{center}%
        \end{minipage}
      }%
    \end{center}%
  \end{figure}%
}

%
% Comando para incluir imágenes en formato apaisado (sin marco).
\newcommand{\figuraApaisadaSinMarco}[5]{%
  \begin{figure}%
    \begin{center}%
    \includegraphics[angle=90,height=#1\textheight,#5]{#2}%
    \caption{#3}%
    \label{#4}%
    \end{center}%
  \end{figure}%
}
% Para las tablas
\newcommand{\otoprule}{\midrule [\heavyrulewidth]}
%
% Nuevo comando para tablas pequeñas (menos de una página).
\newcommand{\tablaSmall}[5]{%
 \begin{table}
  \begin{center}
   \rowcolors {2}{gray!35}{}
   \begin{tabular}{#2}
    \toprule
    #4
    \otoprule
    #5
    \bottomrule
   \end{tabular}
   \caption{#1}
   \label{tabla:#3}
  \end{center}
 \end{table}
}

%
% Nuevo comando para tablas pequeñas (menos de una página).
\newcommand{\tablaSmallSinColores}[5]{%
 \begin{table}[H]
  \begin{center}
   \begin{tabular}{#2}
    \toprule
    #4
    \otoprule
    #5
    \bottomrule
   \end{tabular}
   \caption{#1}
   \label{tabla:#3}
  \end{center}
 \end{table}
}

\newcommand{\tablaApaisadaSmall}[5]{%
\begin{landscape}
  \begin{table}
   \begin{center}
    \rowcolors {2}{gray!35}{}
    \begin{tabular}{#2}
     \toprule
     #4
     \otoprule
     #5
     \bottomrule
    \end{tabular}
    \caption{#1}
    \label{tabla:#3}
   \end{center}
  \end{table}
\end{landscape}
}

%
% Nuevo comando para tablas grandes con cabecera y filas alternas coloreadas en gris.
\newcommand{\tabla}[6]{%
  \begin{center}
    \tablefirsthead{
      \toprule
      #5
      \otoprule
    }
    \tablehead{
      \multicolumn{#3}{l}{\small\sl continúa desde la página anterior}\\
      \toprule
      #5
      \otoprule
    }
    \tabletail{
      \hline
      \multicolumn{#3}{r}{\small\sl continúa en la página siguiente}\\
    }
    \tablelasttail{
      \hline
    }
    \bottomcaption{#1}
    \rowcolors {2}{gray!35}{}
    \begin{xtabular}{#2}
      #6
      \bottomrule
    \end{xtabular}
    \label{tabla:#4}
  \end{center}
}

%
% Nuevo comando para tablas grandes con cabecera.
\newcommand{\tablaSinColores}[6]{%
  \begin{center}
    \tablefirsthead{
      \toprule
      #5
      \otoprule
    }
    \tablehead{
      \multicolumn{#3}{l}{\small\sl continúa desde la página anterior}\\
      \toprule
      #5
      \otoprule
    }
    \tabletail{
      \hline
      \multicolumn{#3}{r}{\small\sl continúa en la página siguiente}\\
    }
    \tablelasttail{
      \hline
    }
    \bottomcaption{#1}
    \begin{xtabular}{#2}
      #6
      \bottomrule
    \end{xtabular}
    \label{tabla:#4}
  \end{center}
}

%
% Nuevo comando para tablas grandes sin cabecera.
\newcommand{\tablaSinCabecera}[5]{%
  \begin{center}
    \tablefirsthead{
      \toprule
    }
    \tablehead{
      \multicolumn{#3}{l}{\small\sl continúa desde la página anterior}\\
      \hline
    }
    \tabletail{
      \hline
      \multicolumn{#3}{r}{\small\sl continúa en la página siguiente}\\
    }
    \tablelasttail{
      \hline
    }
    \bottomcaption{#1}
  \begin{xtabular}{#2}
    #5
   \bottomrule
  \end{xtabular}
  \label{tabla:#4}
  \end{center}
}



\definecolor{cgoLight}{HTML}{EEEEEE}
\definecolor{cgoExtralight}{HTML}{FFFFFF}

%
% Nuevo comando para tablas grandes sin cabecera.
\newcommand{\tablaSinCabeceraConBandas}[5]{%
  \begin{center}
    \tablefirsthead{
      \toprule
    }
    \tablehead{
      \multicolumn{#3}{l}{\small\sl continúa desde la página anterior}\\
      \hline
    }
    \tabletail{
      \hline
      \multicolumn{#3}{r}{\small\sl continúa en la página siguiente}\\
    }
    \tablelasttail{
      \hline
    }
    \bottomcaption{#1}
    \rowcolors[]{1}{cgoExtralight}{cgoLight}

  \begin{xtabular}{#2}
    #5
   \bottomrule
  \end{xtabular}
  \label{tabla:#4}
  \end{center}
}




\graphicspath{ {./img/} }

% Capítulos
\chapterstyle{bianchi}
\newcommand{\capitulo}[2]{
	\setcounter{chapter}{#1}
	\setcounter{section}{0}
	\chapter*{#2}
	\addcontentsline{toc}{chapter}{#2}
	\markboth{#2}{#2}
}

% Apéndices
\renewcommand{\appendixname}{Apéndice}
\renewcommand*\cftappendixname{\appendixname}

\newcommand{\apendice}[1]{
	%\renewcommand{\thechapter}{A}
	\chapter{#1}
}

\renewcommand*\cftappendixname{\appendixname\ }

% Formato de portada
\makeatletter
\usepackage{xcolor}
\newcommand{\tutor}[1]{\def\@tutor{#1}}
\newcommand{\course}[1]{\def\@course{#1}}
\definecolor{cpardoBox}{HTML}{E6E6FF}
\def\maketitle{
  \null
  \thispagestyle{empty}
  % Cabecera ----------------
\noindent\includegraphics[width=\textwidth]{cabecera}\vspace{1cm}%
  \vfill
  % Título proyecto y escudo informática ----------------
  \colorbox{cpardoBox}{%
    \begin{minipage}{.8\textwidth}
      \vspace{.5cm}\Large
      \begin{center}
      \textbf{TFG del Grado en Ingeniería Informática}\vspace{.6cm}\\
      \textbf{\LARGE\@title{}}
      \end{center}
      \vspace{.2cm}
    \end{minipage}

  }%
  \hfill\begin{minipage}{.20\textwidth}
    \includegraphics[width=\textwidth]{escudoInfor}
  \end{minipage}
  \vfill
  % Datos de alumno, curso y tutores ------------------
  \begin{center}%
  {%
    \noindent\LARGE
    Presentado por \@author{}\\ 
    en Universidad de Burgos --- \@date{}\\
    Tutor: \@tutor{}\\
  }%
  \end{center}%
  \null
  \cleardoublepage
  }
\makeatother


% Datos de portada
\title{título del TFG \\Documentación Técnica}
\author{nombre alumno}
\tutor{nombre tutor}
\date{\today}

\begin{document}

\maketitle



\cleardoublepage



%%%%%%%%%%%%%%%%%%%%%%%%%%%%%%%%%%%%%%%%%%%%%%%%%%%%%%%%%%%%%%%%%%%%%%%%%%%%%%%%%%%%%%%%



\frontmatter


\clearpage

% Indices
\tableofcontents

\clearpage

\listoffigures

\clearpage

\listoftables

\clearpage

\mainmatter

\appendix

\apendice{Plan de Proyecto Software}


\section{Introducción}

A continuación se va a describir en profundidad cuál ha sido la planificación del proyecto, semana a semana, organizado en base a sprints.

Para tener una visión apropiada de todo el proceso, se ha usado la extensión ZenHub que se puede añadir a GitHub y que permite observar con más detalle todo lo relacionado con la metodología ágil. El repositorio de este proyecto es \url{https://github.com/mrr0088/Gestor-de-Hemeroteca}, aunque ya se incluyó al usuario ``ubutfgm" para que las observaciones necesarias sean realizadas.

Cabe destacar que, especialmente en los primeros sprints, por falta de experiencia con la herramienta ZenHub, los gráficos de progreso no son representativos del progreso real que se llevó a cabo, debido sobre todo a que los issues que los componen no estaban cerrados correctamente.


\section{Planificación temporal}

Como ya se ha mencionado, el desarrollo se dividió en sprints, equivalentes a una semana, excepto en un sprint de dos semanas. Cada comienzo y final de sprint coincide con la reunión semanal del alumno con los tutores.

En nuestro caso, hemos considerado que los story points asignados a cada tarea simbolicen las horas que se tarda en resolver ese issue. Por ejemplo, si una tarea tiene 5 story points se estima que el alumno la realizará en 5 horas.

\textbf{Primer Sprint (2 de Febrero - 10 de Febrero):} 

Este tramo inicial fue dedicado a realizar investigación sobre las herramientas que se usarían para desarrollar el producto final. Se buscó información sobre 3 elementos o partes distintas:

\begin{itemize}

\item Herramienta para desarrollar la GUI: Concretamente, se debatió acerca de si sería más adecuado crear una aplicación web o una aplicación interna. Para tomar una decisión, se tuvieron en cuenta 4 herramientas. Por la parte de aplicación web, las candidatas fueron Jupyter y Flask; y por la parte de aplicación de escritorio, PyQt y Kivy. Tras practicar con ambas herramientas, se tomó la decisión de programar la GUI en un notebook de Jupyter para posteriormente publicarlo en un Dashboard.

\item Base de datos para almacenar la información: Como se debía almacenar numerosa y variada información sobre las noticias elegidas, se necesitaba de una estructura donde almacenar todos esos datos. La decisión más relevante en este apartado era si usar una base de datos relacional o no relacional, o dicho de otra manera, SQL o NoSQL. Las alternativas fueron AlchemySql (relacional), PyLucene (no relacional) y PyMongo (no relacional).

\item Herramienta para escribir la documentación: Aquí sólo se plantearon dos alternativas inicialmente: OpenOffice o Latex. Por cuestiones de formato y de preferencia de los miembros del proyecto, se optó por Latex, aprovechando la plantilla para Trabajos de Fin de Grado que ya había sido creada y que estaba a disposición de los alumnos.

\end{itemize}

\imagen {Issues1}{Issues del primer sprint}

\imagen {Report1}{Progreso de resolución de tareas del primer sprint}

\textbf{Segundo Sprint (10 de Febrero - 16 de Febrero):}

La prioridad de este sprint era finalizar la investigación y búsqueda de herramientas para desarrollar la aplicación, que concluyó con la elección de Pymongo como base de datos, y los notebook de Jupyter como herramienta de interfaz.
Una vez elegidas las diferentes herramientas necesarias para el desarrollo del proyecto, el siguiente paso consistió en crear un pequeño programa que probara la librería PDFMiner y que contuviera algo de funcionalidad básica como contar la frecuencia de una palabra dentro de un documento.
Aprovechando este script, se realizaron pruebas con un snippet llamado fileupload, que permite cargar ficheros en el código usando una GUI clásica en la que puedes navegar por el equipo buscando la ubicación del mencionado fichero. El resultado final era un programa que  permite al usuario buscar un fichero en el equipo y, si este archivo tiene una extensión pdf, se permitirá buscar la frecuencia de las palabras que introduzca el propio usuario.

\imagen {Issues2}{Issues del segundo sprint}

\imagen {Report2}{Progreso de resolución de tareas del segundo sprint}

\textbf{Tercer Sprint (16 de Febrero - 23 de Febrero):}

Habiendo finalizado la elección de todas las herramientas necesarias, y probado el funcionamiento básico para lectura de PDF subiendo archivos a mano, la siguiente prioridad fue crear un primer prototipo de web scrapping.
Para lograr esto, se dividió la funcionalidad en dos etapas:
\begin{itemize}

\item Un primer script que realiza la lectura web de noticias en diferentes medios (para este prototipo, se leían noticias sólo del diario “Público” y, más en concreto, de la sección de Igualdad), almacenando en la base de datos estructuras formadas por el titular, autor, fecha, texto de la noticia, procedencia y un link a la misma.
Para recoger el texto html de las páginas web de los diarios, se usaron las librerías de url2lib y BeautifulSoup; la primera para hacer una llamada a la url indicada y obtener todo (o casi todo, en las próximas etapas explicaremos qué faltaba) el texto HTML de la página de cada noticia, y la segunda para poder manipular todo el HTML recibido previamente de forma más sencilla sin tener que hacer una división manual por etiquetas, simplemente indicamos al contenido de qué etiquetas queremos acceder.
 
\item Un segundo script que, en función de una palabra clave introducida por el propio usuario, consulte en la base de datos y extraiga todas las noticias que contengan esa palabra clave, en un formato de tabla HTML y mostrando la información más importante de la noticia junto con el número de apariciones que hace esa palabra en el texto de las mismas.
Para poder mostrar esto en la salida del notebook, necesitamos las librerías de pandas, ipywidgets y qgrid. Con la librería de pandas podemos organizar la información devuelta en la consulta de la base de datos a un dataframe, que puede ser convertido a una tabla html. Con los widgets, podemos mostrar controles básicos como botones, entradas de texto, listas desplegables, etc… sin necesidad de escribir a mano el código HTML. En esta versión, se usó la librería qgrid para transformar el dataframe en una tabla interactiva, que permitía reordenar el orden de las filas dando prioridad a unas columnas u otras. No obstante, no permitían formatear las celdas por separado, lo cual nos hizo descartarlo para posteriores versiones.

\end{itemize}

\imagen {Issues3}{Issues del tercer sprint}

\imagen {Report3}{Progreso de resolución de tareas del tercer sprint}

\textbf{Cuarto Sprint (23 de Febrero - 2 de Marzo):}

Habiendo realizado una primera aproximación tanto a la lectura de PDF como de web, este sprint se centró en mejorar tres apartados:

\begin{itemize}

\item En la parte de PDF, se mejoró el script, aunque a costa de retirar temporalmente la librería que permitía seleccionar un fichero de forma manual, para que recorriera y almacenara todos los archivos PDF, tanto del directorio que se pasaba como argumento como de todos los que colgaban de él. Una vez recorridos todos, se guardan las noticias en la base de datos, de forma similar a como se hacía con las noticias en la web pero cambiando el atributo de la procedencia.

\item En la parte de web scrapping, se actualizaron dos aspectos importantes:

\begin{itemize}

\item Por una parte, usando Selenium además de BeautifulSoup se consiguió obtener los comentarios de las noticias, que en el caso de “Público”, se generaban con una API de terceros que impedía acceder a su HTML de forma normal como se hacía con el resto de la noticia. El método era algo pesado, ya que selenium necesita abrir en el ordenador que se esté ejecutando las pestañas a las que accede, pero como ya se ha mencionado, fue la alternativa que se encontró más idónea para poder acceder al HTML de los comentarios de las noticias.

\item Por otro lado, se cambió el formato de la tabla, dejando de usar qgrid para el DataFrame y mostrándola en un HTML normal, al que luego se aplicaba estilo manualmente usando un script de Javascript. Ésto es posible ya que uno de los widgets de Ipython/Jupyter permite introducir código HTML de forma manual, y es la vía que también utilizamos para introducir código JavaScript dentro del output del notebook. Con este cambio, tenemos flexibilidad total para manipular las celdas de la tabla para, por ejemplo, formatear como hiperenlaces las url de las noticias que se muestran en la tabla.


\end{itemize}

\end{itemize}

\imagen {Issues4}{Issues del cuarto sprint}

\imagen {Report4}{Progreso de resolución de tareas del cuarto sprint}

\textbf{Quinto Sprint (2 de Marzo - 10 de Marzo):}

En esta etapa se decidió que en el muestreo de resultados, sería interesante incluir gráficos que acompañaran al resto de la información, para poder mostrar la mayor cantidad de datos posibles de diversas formas. Por tanto, una de las tareas fue elegir una librería de gráficos adecuada a nuestro contexto. 
Las candidatas iniciales fueron Bokeh y SeaBorn. Se probó Bokeh y, debido a unos problemas relacionados con la generación de archivos HTML, de decidió descartar. Por falta de horas durante esa semana, se trasladaron las pruebas con SeaBorn al siguiente sprint.
La otra gran tarea a la que se dedicó tiempo durante este sprint fue a aumentar los medios a los que se realizaba web scrapping, creando un script para recorrer y almacenar las noticias de ElDiario. En este caso, los comentarios no se generaban a partir de una API de terceros, pero debido a la complejidad del código HTML obtenido, y visto necesario el uso de Selenium una vez más, se decidió separar la obtención de comentarios y aplazar esa parte al siguiente sprint.

\imagen {Issues5}{Issues del quinto sprint}

\imagen {Report5}{Progreso de resolución de tareas del quinto sprint}

\textbf{Sexto Sprint (10 de Marzo - 16 de Marzo):}

La mayor parte de esta semana se dedicó a terminar tareas de la semana anterior que habían sido subestimadas en horas, y separadas en tareas entre este sprint y el anterior.
En el apartado de los gráficos, se realizaron pruebas con la otra librería candidata, Seaborn, y por su comodidad y sencillez para mostrar gráficos simples (Como gráficos de barras para evaluar las menciones a una palabra por autor o día de la semana), se decidió que por el momento sería la alternativa elegida para acompañar en la interfaz.
Respecto al web scrapping de ElDiario, se implementó el código necesario junto con el uso de Selenium adecuado para poder obtener los comentarios de las noticias de ese medio, que era la funcionalidad que no dio tiempo a terminar en la anterior semana.
Aparte de estas tareas, se dedicó cierta cantidad de tiempo a investigar formas de publicar esta aplicación web en algún servidor. Para comenzar, se realizó cierta investigación sobre Docker, además de ver vídeos de la PyCon de Londres 2016, donde se trataba el tema de subir notebook a un servidor y poder visualizarlos como un dashboard. Por problemas de límite de tiempo, no se pudo avanzar mucho más este aspecto del proyecto.

\imagen {Issues6}{Issues del sexto sprint}

\imagen {Report6}{Progreso de resolución de tareas del sexto sprint}

\textbf{Séptimo Sprint (16 de Marzo - 23 de Marzo):}

En este sprint se dedicó parte del tiempo a solucionar un problema recurrente en el que se mostraban sucesivamente todas las tablas de resultados de todas las palabras buscadas, provocando que la interfaz resultante se llenara de tablas innecesarias. De la misma manera, con los gráficos de Seaborn, que usa MatplotLib para funcionar, sucedía algo similar, provocando que todos los gráficos generados se fueran acumulando en una misma zona de dibujado que comparten por requisitos de la librería.
Para solucionar esto hicieron falta dos medidas distintas. Para la parte de las tablas html, el problema se solucionó desde el código de Python, cambiando el planteamiento de la función que se ejecutaba al pulsar en el botón de búsqueda, y haciendo globales algunas variables necesarias para la generación de estas tablas. Para la parte de los gráficos, como este problema se estaba produciendo por limitaciones de la librería, hubo que buscar una solución más “rudimentaria” y controlar, por medio de una función JavaScript (introducida en un widget HTML), la visualización de la tabla requerida y la no visualización del resto de tablas obsoletas.
El resto del tiempo de esa semana se dedicó a la investigación de la herramienta TextBlob, una opción muy interesante para realizar minería de datos de texto. En términos generales, es una herramienta que, a partir de los textos recibidos, puede traducirles, detectar idiomas, corregir palabras, dividir párrafos en frases, y frases en palabras, devolver etiquetas de las frases o palabras… etc. Aparte de toda esta funcionalidad, muy útil para casos como el de este proyecto, en el que se manejan textos, tiene sus propios clasificadores de datos. 
Para usar esta herramienta, se necesita descargar el NLTK (Natural Language ToolKit), que contiene todos los conjuntos de datos necesarios para poder llevar a cabo estas manipulaciones de texto y clasificaciones de los mismos. 
Debido a problemas con la descarga y tamaño de los archivos, se cerró la tarea en ese punto, y se retomó en el siguiente sprint.
Al término de esta semana, se realizó una reunión con el cliente para enseñar los progresos alcanzados hasta ese punto y tomar decisiones acerca de qué pasos tomar a continuación. 

\imagen {Issues7}{Issues del séptimo sprint}

\imagen {Report7}{Progreso de resolución de tareas del séptimo sprint}

\textbf{Octavo Sprint (23 de Marzo - 30 de Marzo):}

La tarea más importante realizada esa semana fue la de hacer otra aproximación al web scrapping basada en obtener los datos a partir de los RSS de los medios en los que, hasta el momento, se estaban leyendo desde las páginas web convencionales. El resultado fue un pequeño programa que te permitía seleccionar un medio de una lista de diarios y una palabra clave y, en tiempo real, realizar el web scrapping del RSS de ese medio y almacenarlo en la base de datos a la vez que se mostraba en una tabla los resultados obtenidos. Para ello se usó la librería feedparser.
Este acercamiento era mucho más rápido que el web scrapping convencional, pero a cambio sacrifica la posibilidad de poder acceder a los comentarios de las noticias, y los casos de algunos medios, no todo el texto de la noticia es accesible a través del RSS.
Como ha sido mencionado en el resumen del sprint anterior, debido a la falta de capacidad en la máquina virtual en la que se estaba trabajando, no se podía descargar todo el conjunto de datos de NLTK de forma correcta, lo cual se intentó solucionar inicialmente clonando la máquina virtual a una de almacenamiento dinámico y editando la capacidad de almacenamiento de las particiones de la misma. 
Esta solución no tuvo resultados positivos, por tanto, se reinstaló desde el principio una nueva máquina virtual que evitara estos problemas y que además corrigiera la instalación de algunas librerías que habían producido problemas previamente.
Por último, se probó una nueva librería de gráficos que sustituyera a Seaborn, denominada Bqplot, la cual se asemejaba en funcionalidad a Bokeh, pero evitando esa necesidad de archivos HTML que nos habían llevado a descartar Bokeh en un primer momento. Las pruebas fueron muy positivas en gráficos de barras y lineales.

\imagen {Issues8}{Issues del octavo sprint}

\imagen {Report8}{Progreso de resolución de tareas del octavo sprint}

\textbf{Noveno Sprint (31 de Marzo - 13 de Abril):}

La duración de este sprint fue el doble de lo normal (dos semanas en vez de una) debido a que coincidió con las vacaciones de Semana Santa. Por tanto, el día que debía haberse realizado la reunión habitual no era lectivo, y hubo que posponerlo. 
Por otro lado, como había suficientes tareas/horas de trabajo como para dividirlo en más de un sprint, se dio por finalizado el sprint a los 14 días aunque la siguiente reunión fuera después de 21.
En lo relacionado a tareas relacionadas con sprints anteriores, se consiguió solucionar el problema por el cual no se aplicaban los scripts de JavaScript, aunque para ello hubo que buscar una forma alternativa de implementarles.
Además, se consiguió finalizar con éxito la instalación y pruebas con TextBlob, llevando a cabo la instalación del NLTK de una manera alternativa que ocupaba mucho menos espacio en memoria. Una vez instalado y probado, se hizo un sencillo prototipo con un conjunto de datos de prueba minúsculo de diferenciador de frases machistas y feministas, tanto con TextBlob como con SciKitLearn, aprovechando código de scripts ya existentes.
También se mejoraron las operaciones de almacenamiento en base de datos, evitando introducir noticias duplicadas, para que el usuario final no se tenga que preocupar de filtrar las que estén duplicadas de las que no.
En lo respectivo a la publicación del notebook como aplicación web, se realizaron instalaciones de varias herramientas de Jupyter Dashboards y de Docker, pero debido a impedimentos de configuración de la BIOS y del anfitrión, no se puedo avanzar más en la máquina virtual, y se decidió repetir el proceso en el anfitrión en el siguiente sprint.
Por último, se añadió el diario “El Mundo” a los disponibles en el web scrapping por RSS, para ir aumentando el rango ideológico de las alternativas disponibles.

\imagen {Issues9}{Issues del noveno sprint}

\imagen {Report9}{Progreso de resolución de tareas del noveno sprint}

\textbf{Décimo Sprint (14 de Abril - 21 de Abril):}

En este periodo se continuó la investigación y las pruebas de despliegue de aplicación de los notebook en una máquina de Docker, aunque en esta ocasión se realizaron en el Sistema Operativo anfitrión (Windows) en vez de en la máquina virtual, para comprobar si al cambiar el entorno conseguíamos solucionar los problemas técnicos sufridos durante la configuración de Docker.

Se realizaron además las primeras pruebas de clasificación de textos con un dataset real, compuesto por un grupo de noticias de diferentes medios de comunicación tratando el debate de los vientres de alquiler, etiquetadas a partes iguales como "A Favor" y "En Contra" del tema que hablaban. Las pruebas fueron realizadas tanto en TextBlob, dividiéndolas en sentencias, como en SciKit-Learn, dividiéndolas en palabras.

También se dedicó parte de la semana a la documentación del trabajo realizado hasta la fecha.

\imagen {Issues10}{Issues del décimo sprint}

\imagen {Report10}{Progreso de resolución de tareas del décimo sprint}

\textbf{Undécimo Sprint (22 de Abril - 27 de Abril):}

Tras los fracasos en la configuración sufridos con Docker, se investigó la alternativa de Cloud Foundry, plataforma open source basada en contenedores que da soporte a la subida de aplicaciones a la nube.

En lo respectivo a la clasificación del Dataset real, se continuaron las pruebas, en este caso probando distintos permutadores (Leave One Out, Leave P Out y ShuffleSplit), para hacer validación cruzada y separar las noticias que teníamos en conjuntos de entrenamiento y conjunto de prueba. 

Además, se investigó la librería de Lime, que permite crear explicaciones detalladas de las predicciones, y se realizaron los primeros experimentos con la misma.

\imagen {Issues11}{Issues del undécimo sprint}

\imagen {Report11}{Progreso de resolución de tareas del undécimo sprint}

\textbf{Duodécimo Sprint (28 de Abril - 4 de Mayo):}

En este punto, se decidió dar más prioridad a crear una aplicación con funcionalidad completa que a seguir invirtiendo más tiempo en intentar configurar un contenedor para poder subir los notebooks de Jupyter con formato de aplicación web. 

Por ello, se comenzó el desarrolló de una aplicación con Flask que reuniera toda la funcionalidad implementada hasta el momento, que estaba repartida en varios archivos sin relación entre ellos, con objetivo de desplegarla en modo local en lugar de subirla a un servidor. 

Se creó un esqueleto con la funcionalidad más básica de lectura de PDF, Web Scraping y persistencia en una base de datos NoSQL, sobre la que construir en las siguientes semanas.

\imagen {Issues12}{Issues del duodécimo sprint}

\imagen {Report12}{Progreso de resolución de tareas del duodécimo sprint}

\textbf{Decimotercer Sprint (5 de Mayo - 11 de Mayo):}

Durante esta semana se continuó la construcción de la aplicación en Flask, como la implementación de las inserciones desde la web a la base de datos, el traslado del código de clasificación de texto de Jupyter a Flask y una mejora en la lectura de PDF.

Aparte de esto, se dedicó un tiempo a solucionar problemas más pequeños como correciones en el uso de rutas literales, sustituir el uso de librerías obsoletas por sus alternativas vigentes, la inclusión de un fichero .txt con las librerías necesarias para ejecutar la aplicación, y completar la experimentación de Lime con los ejemplos de la documentación.

\imagen {Issues13}{Issues del decimotercer sprint}

\imagen {Report13}{Progreso de resolución de tareas del decimotercer sprint}

\textbf{Decimocuarto Sprint (12 de Mayo - 18 de Mayo):}

Este sprint se centró en avanzar todo lo posible la funcionalidad relacionada con la clasificación de texto. Concretamente, se consiguió desde la librería de NLTK una lista de palabras irrelevantes en español a tener en uenta para evitar ruido en los entrenamientos y predicciones. 

En lo respectivo a la aplicación web, también se centraron esfuerzos en la vista de resultados de noticias y las posibilidades de etiquetado de las mismas, permitiendo que las noticias tuvieran varios datasets, que se pudiera añadir un dataset entero mediante un documento XML y se hizo una primera aproximación a un etiquetado manual, pero que sería sustituido por una alternativa más coherente en la siguiente semana.

\imagen {Issues14}{Issues del decimocuarto sprint}

\imagen {Report14}{Progreso de resolución de tareas del decimocuarto sprint}

\textbf{Decimoquinto Sprint (19 de Mayo - 25 de Mayo):}

En esta semana se completó casi toda la funcionalidad restante en la tabla de resultados de noticias, incluyendo una barra de progreso que indicaba las probabilidades que tenía una noticia de pertenecer a una clase o a otra, la posibilidad de que el usuario realice un etiquetado manual a cada noticia, y la posibilidad de ver la explicación de la predicción en una pestaña aparte, acompañada de un gráfico para comprender mejor los resultados.

También se dedicó parte del tiempo a mejorar la coherencia de urls para separar distintas funcionalidades en distintas pestañas, para evitar sobrecargar de complejidad las peticiones de los formularios, sobre todo en las pestañas de lectura de PDF y XML.

\imagen {Issues15}{Issues del decimoquinto sprint}

\imagen {Report15}{Progreso de resolución de tareas del decimoquinto sprint}

\textbf{Decimosexto Sprint (26 de Mayo - 2 de Junio):}

Teniendo gran parte de la funcionalidad implementada, en este Sprint se centró en mejorar el estilo gráfico de la interfaz de la aplicación. Para ello, se usó Bootstrap, primero experimentando con la versión para Flask y posteriormente se decidió permanecer con el modo de uso que tendría en un proyecto web normal, importando los archivos CSS y JS correspondientes. Aprovechando que había que utilizar JQuery para poder aplicar Bootstrap, se refactorizó el código existente de JavaScript en JQuery, reduciendo en un gran tamaño el mismo.

Además del estilo, también se actualizó el modo de lectura de los PDF, creando una estructura dentro de los archivos simple que facilitara la posterior lectura e inserción en base de datos, obligando a todos los documentos PDF que sigan la misma estructura para que puedan ser leídos. 

También se modificó la vista de resultados de la explicación de Lime, aprovechando su característica de poder mostrar una vista de resultados propia exportada como HTML, sin tener que recurrir a librerías de gráficos de terceros.

\imagen {Issues16}{Issues del decimosexto sprint}

\imagen {Report16}{Progreso de resolución de tareas del decimosexto sprint}

\textbf{Decimoséptimo Sprint (2 de Junio - 9 de Junio):}

A estas alturas ya se estaban realizando las últimas mejoras en la funcionalidad, como una actualización de la lectura de los RSS de las páginas web, de los XML y de los PDF. Aparte de ésto, se terminaron los capítulos de la memoria dedicado a los conceptos teóricos y objetivos del proyecto.

\imagen {Issues17}{Issues del decimoseptimo sprint}

\imagen {Report17}{Progreso de resolución de tareas del decimoseptimo sprint}

\textbf{Decimoctavo Sprint (9 de Junio - 15 de Junio):}

En lo relacionado con la documentación, se completaron los capítulos de aspectos relevantes, herramientas utilizadas y las conclusiones y líneas futuras del proyecto.

A lo largo de esta semana también se finalizó la funcionalidad, añadiendo una pestaña para ver gráficos de las noticias almacenadas, además de arreglar los problemas que surgieron al añadir la posibilidad de añadir datasets adicionales. Por último, se implementó la paginación en la lista de resultados, para partir las noticias en pequeños grupos y evitar tiempos de carga demasiado largos.

\imagen {Issues18}{Issues del decimoctavo sprint}

\imagen {Report18}{Progreso de resolución de tareas del decimoctavo sprint}

\textbf{Decimonoveno Sprint (16 de Junio - 22 de Junio):}

Habiendo finalizado la funcionalidad, los objetivos de esta semana se enfocaron en realizar la documentación de los anexos y terminar la de la memoria, a la vez que se realizaban pruebas sobre el código y se realizó refactorización sobre el mismo.

También se arreglaron algunos problemas relacionados con el almacenamiento de nuevas noticias.

\imagen {Issues19}{Issues del decimonoveno sprint}

\imagen {Report19}{Progreso de resolución de tareas del decimonoveno sprint}

\section{Estudio de viabilidad}

\subsection{Viabilidad económica}

Para poder hacer un uso correcto de la aplicación, no hay que hacer ninguna inversión económica en software de ningún tipo, ya que todas las librerías que se han usado son open-source.

Más allá de lo comprendido por la propia herramienta, el único coste imprescindible sería el del propio equipo en el que se instale y configure. Aunque algunas funcionalidades tengan tiempos de carga más notables que otras, no debería ser impedimento para que la aplicación funcionara correctamente en equipos de diferentes prestaciones. Respecto al Sistema Operativo, como la aplicación se ejecuta en sistemas Linux, no es necesario ni realizar la compra de un Sistema Operativo de pago.

Cabe destacar que en el contexto de este proyecto la aplicación no está publicada en un servidor web, pero que en un futuro podría llevarse a cabo esta operación. De ser así, habría que evaluar diferentes alternativas de host compatibles y comprobar si alguna de estas alternativas son de pago o gratuitas.

\subsubsection{Viabilidad económica del desarrollo}

En lo relacionado con costes de personal, hay varios datos de los que se podrían obtener algunos presupuestos hipotéticos.

Por ejemplo, según la información proporcionada por \emph{ZenHub} en el repositorio del proyecto, se estima que se han invertido 556 horas en el desarrollo, a lo largo de los 6 meses del mismo.

Si se propusiera un sueldo para el programador de 10\euro{} la hora, tendríamos el siguiente cálculo:

\[10  /\textsf{hora} * 556 \textsf{horas} = 5560 \]

De esta cantidad inicial, habría que reflexionar acerca de cuánto tiempo ha sido dedicado a programar la aplicación final, y cuánto tiempo fue dedicado a la investigación y programación de funcionalidades que no acabaron incluidas en el producto, ya que podría considerarse que esas horas no forman parte del desarrollo de la herramienta, y por tanto ser descontadas de la cantidad calculada.

\subsection{Viabilidad legal}

A continuación, se expondrán de forma breve las licencias de todas las herramientas y librerías utilizadas en la creación del producto final.

\begin{table}[]
	\centering
	\caption{Tabla de bibliotecas y sus licencias}
	\label{tabla:Licencias}
	\rowcolors {2}{gray!35}{}
	\begin{tabular}{l l l l}
	\toprule
	Librería     & Versión & Descripción                                                     & Licencia                \\ 	\midrule
	Flask       & 0.12.2     &  \begin{tabular}[c]{@{}l@{}}Biblioteca para crear\\ aplicaciones web en Python. \end{tabular}               & BSD    \\
	Jinja2       & 2.8     &  \begin{tabular}[c]{@{}l@{}} Motor para integrar Python \\ documentos HTML. \end{tabular}               & BSD                     \\ 	
	Pymongo        & 3.4.0  & \begin{tabular}[c]{@{}l@{}}Extensión de  MongoDB (Base\\ de datos NoSQL) para Python\\  \end{tabular} & Apache license v2.0                     \\ 	
	Scikit-Learn & 0.18  & \begin{tabular}[c]{@{}l@{}}Biblioteca para aprendizaje\\ automático en Python.\end{tabular}	               & BSD                     \\ 	
	Lime     & (Desconocida)    & \begin{tabular}[c]{@{}l@{}}biblioteca Python para crear\\ explicaciones de clasificaciones.\end{tabular}          & BSD                     2-clause \\ 	
	Beautiful Soup       & 4.6.0   & \begin{tabular}[c]{@{}l@{}}Biblioteca para hacer web\\ scraping en Python.\end{tabular} & MIT \\ 
PDFMiner         & (Desconocida)   &\begin{tabular}[c]{@{}l@{}} Biblioteca para lectura\\ de PDF en Python.                               \end{tabular} &  OSI Approved          \\ 
	Feedparser   & 5.2.1   & \begin{tabular}[c]{@{}l@{}}Biblioteca Python para hacer\\ lectura de RSS.\end{tabular}                      & OSI Approved              \\ 
	Bootstrap    & 3.3.7    & \begin{tabular}[c]{@{}l@{}}Biblioteca para desarrollar\\ aplicaciones web responsive.\end{tabular}                      & MIT  	\\
	JQuery    & 3.2.1    & \begin{tabular}[c]{@{}l@{}}Biblioteca para optimizar\\ JavaScript.\end{tabular}
& MIT                 \\
	Chart.js    & 2.0    & \begin{tabular}[c]{@{}l@{}}Biblioteca para crear\\ gráficos en Javascript y HTMl.\end{tabular}                     
	& MIT \\ \bottomrule
\end{tabular}
\end{table}



\apendice{Especificación de Requisitos}

\section{Introducción}

En esta sección se describirán las funciones para las que la herramienta fue desarrollada. Se hará una detallada definición tanto de objetivos como de sus requisitos (funcionales y no funcionales) y también se hará un estudio en profundidad de los casos de uso.

\section{Objetivos generales}

Como ya se describió previamente en a memoria, hay un conjunto de objetivos en torno a los cuales se estructura el desarrollo del proyecto. A continuación se volverán a resumir los relacionados con la funcionalidad y las técnicas empleadas:

\begin{itemize}

\item El principal objetivo es crear una aplicación web en la que el usuario pueda almacenar noticias tanto por medio de web scraping como lectura de PDF.

\item Se busca que el usuario pueda añadir etiquetas  relacionadas con diferentes temas de debate a estas noticias, tanto manualmente como guardando una predicción realizada con aprendizaje automático.

\item Otro objetivo es que el usuario pueda añadir temas de debate nuevos que contengan dos clases, para luego poder etiquetar estas noticias en función de estas nuevas clases.

\item En caso de que el usuario ya tuviera un corpus preparado antes de comenzar a usar la aplicación, tiene la posibilidad de añadir un corpus mediante un archivo XML con todas las noticias que lo forman y sus etiquetas correspondientes.

\item La aplicación debe ser responsive en los diferentes equipos en los que se use, y fácil de usar para nuevos usuarios.


\end{itemize}

\section{Catalogo de requisitos}

\subsection{Requisitos Funcionales}

Estos son los requisitos funcionales que la herramienta debe cumplir:

\begin{itemize}

\item RF - 1: Los usuarios deben poder visualizar un conjunto de noticias a partir de la búsqueda por palabra clave.

\item RF - 2: Los usuarios podrán etiquetar manualmente las noticias en función del tema de debate seleccionado.
\begin{itemize}
\item RF - 2.1: Los usuarios podrán elegir de una lista de opciones el tema de debate mostrado en ese momento.
\end{itemize}

\item RF - 3: Los usuarios pueden ver una explicación gráfica para la predicción de la etiqueta de una noticia.
\begin{itemize}
\item RF - 3.1: Los usuarios también podrán guardar esa predicción como la etiqueta final de la noticia.
\end{itemize}

\item RF - 4: Los usuarios podrán guardar las noticias guardadas en archivos PDF almacenados en una carpeta concreta del proyecto.

\item RF - 5: Los usuarios podrán guardar nuevas noticias de 
los RSS de varios medios de comunicación españoles.
\begin{itemize}
\item RF - 5.1: Los usuarios podrán elegir de una lista de opciones el medio de comunicación del que ver las nuevas noticias.
\end{itemize}

\item RF - 6: Los usuarios podrán añadir un corpus de noticias entero, indicando el tema de debate y la etiqueta de cada noticia, a través de un archivo XML.

\item RF - 7: Los usuarios podrán crear nuevos temas de debate, escribiendo el nombre del propio tema y el nombre de dos clases, que se usarán posteriormente de etiquetas para ese tema.

\end{itemize}

\subsection{Requisitos No Funcionales}

Por otro lado, estos son los requisitos no relacionados con la funcionalidad que se deben tener en cuenta:

\begin{itemize}

\item RNF - 1: La interfaz de las pestañas no debe impedir su correcto funcionamiento en pantallas de diferentes tamaños.

\item RNF - 2: A pesar de que se produzcan picos de carga de trabajo, el funcionamiento general no debe producir tiempos de carga excesivos.

\item RNF - 3: La herramienta debe poder ejecutarse en cualquier navegador (siempre que soporten librerías como Bootstrap o JQuery).

\end{itemize}

\section{Especificación de requisitos}

A continuación, se describirán con detalle los casos de uso de los que se compone el proyecto.

\subsection{Diagrama de casos de uso}

\imagen{casosdeuso}{Diagrama de casos de uso}

\subsection{Descripción de casos de uso}

\clearpage
\subsubsection{CU1 - Consulta de Noticias}

\begin{center}
\begin{tabular}{ | m{3cm} | m{10cm}| } 
\hline
Caso de uso & Consulta de Noticias \\ 
\hline
Autor & Miguel Rodríguez \\ 
\hline
Versión & 1.0 \\ 
\hline
Requisitos & RF - 1: Los usuarios deben poder visualizar un conjunto de noticias a partir de la búsqueda por palabra clave \\ 
\hline
Descripción & El usuario introduce una palabra clave para la búsqueda de las noticias que la contengan \\
\hline
Precondiciones & Las noticias tienen que estar almacenadas en la base de datos. 

La palabra debe aparecer en la noticia para que sea devuelta. \\
\hline
Acciones & En la pestaña inicial, el usuario introduce en el cuadro de texto la palabra clave con la que quiere realizar la búsqueda y pulsa el botón ``Buscar'' \\
\hline
Postcondiciones & Se cargará una tabla con las noticias que contenían la palabra \\
\hline
Excepciones & En caso de que la palabra no se encuentre en ninguna noticia, no se devolverán filas de la tabla de resultados \\
\hline 
Importancia & Alta \\
\hline 
\end{tabular}
\captionof{table}{Caso de uso 1}
\end{center}
\clearpage
\subsubsection{CU2 - Etiquetado Manual}

\begin{center}
\begin{tabular}{ | m{3cm} | m{10cm}| } 
\hline
Caso de uso & Etiquetado Manual \\
\hline
Versión & 1.0 \\  
\hline
Autor & Miguel Rodríguez \\ 
\hline
Requisitos & RF - 2: Los usuarios pueden etiquetar manualmente las noticias en función del tema de debate seleccionado

RF - 2.1: Los usuarios podrán elegir de una lista de opciones el tema de debate mostrado en ese momento \\ 
\hline
Descripción & El usuario selecciona un tema de debate y una de las dos clases para la noticia que quiera etiquetar \\
\hline
Precondiciones & La noticia que se va a etiquetar tiene que estar en la tabla de resultados. El tema de debate sobre el que se quiera etiquetar tiene que existir previamente. \\
\hline
Acciones & El usuario selecciona de una lista desplegable el tema de debate en el que quiera etiquetar las noticias. Una vez seleccionado, pulsando el botón de ``Etiquetar'' de cada fila podrá marcar cuál de las dos etiquetas quiere aplicar a esa noticia. Pulsando en ``Guardar Cambios'', se actualizarán todas las etiquetas modificadas. \\
\hline
Postcondiciones & Se añadirá o etiquetará esa etiqueta a esa noticia. Esa noticia pasa a formar parte del corpus de entrenamiento para ese tema de debate. \\
\hline
Excepciones & Si se realizan etiquetados pero no se guardan los resultados, no se aplicarán los cambios \\
\hline 
Importancia & Alta \\
\hline 
\end{tabular}
\captionof{table}{Caso de uso 2}
\end{center}

\clearpage
\subsubsection{CU3 - Creación de Explicación}

\begin{center}
\begin{tabular}{ | m{3cm} | m{10cm}| } 
\hline
Caso de uso & Creación de explicación \\ 
\hline
Versión & 1.0 \\ 
\hline
Autor & Miguel Rodríguez \\ 
\hline
Requisitos & RF - 3: Los usuarios pueden ver una explicación gráfica para la predicción de la etiqueta de una noticia.

RF - 3.1: Los usuarios también podrán guardar esa predicción como la etiqueta final de la noticia.\\ 
\hline
Descripción & El usuario selecciona la opción de explicación para una noticia de la tabla de resultados. \\
\hline
Precondiciones & La noticia que se va a etiquetar tiene que estar en la tabla de resultados. Debe existir un corpus de entrenamiento para el tema de debate sobre el que se va a realizar la predicción. \\
\hline
Acciones & El usuario selecciona de una lista desplegable el tema de debate en el que quiera etiquetar las noticias. Una vez seleccionado, pulsando el botón de ``Clasificar'' de una noticia en concreto cargará una página con la explicación correspondiente. Pulsando en ``Guardar Predicción'', se guardará el resultado de la predicción como etiqueta. \\
\hline
Postcondiciones & Se etiquetará la noticia con el valor de la clase ganadora en la predicción \\
\hline
Excepciones & Si no hay un corpus de entrenamiento, la explicación no se cargará correctamente \\
\hline 
Importancia & Alta \\
\hline 
\end{tabular}
\captionof{table}{Caso de uso 3}
\end{center}

\clearpage
\subsubsection{CU4 - Lectura de PDF}

\begin{center}
\begin{tabular}{ | m{3cm} | m{10cm}| } 
\hline
Caso de uso & Lectura de PDF \\ 
\hline
Versión & 1.0 \\ 
\hline
Autor & Miguel Rodríguez \\ 
\hline
Requisitos & RF - 4: Los usuarios podrán guardar las noticias guardadas en archivos PDF almacenados en una carpeta concreta del proyecto\\ 
\hline
Descripción & Se muestra una lista de noticias leídas en archivos PDF y el usuario escoge si guardar las nuevas o no. \\
\hline
Precondiciones & Debe haber archivos PDF que contengan noticias en el directorio adecuado. 

Además, deben de seguir el formato especificado para que se puedan leer correctamente. \\
\hline
Acciones & En la pestaña de lectura de PDF, el usuario pulsará el botón de leer noticias. 

Cuando se hayan leído y se muestre la lista, el usuario puede pulsar el botón de ``Guardar Nuevas'' para almacenar las noticias no previamente guardadas. \\
\hline
Postcondiciones & Se almacenarán las noticias no guardadas previamente en la base de datos. \\
\hline
Excepciones & Si no hay archivos PDF no se mostrará nada en la tabla de resultados. 

Si no tienen el formato indicado, no se podrá leer esas noticias. \\
\hline 
Importancia & Media \\
\hline 
\end{tabular}
\captionof{table}{Caso de uso 4}
\end{center}

\clearpage
\subsubsection{CU5 - Lectura de RSS}

\begin{center}
\begin{tabular}{ | m{3cm} | m{10cm}| } 
\hline
Caso de uso & Lectura de RSS \\ 
\hline
Versión & 1.0 \\ 
\hline
Autor & Miguel Rodríguez \\ 
\hline
Requisitos & RF - 5: Los usuarios podrán guardar nuevas noticias de 
los RSS de varios medios de comunicación españoles.

RF - 5.1: Los usuarios podrán elegir de una lista de opciones el medio de comunicación del que ver las nuevas noticias.
\\ 
\hline
Descripción & El usuario elige un medio de comunicación y después puede guardar las nuevas noticias de ese medio. \\
\hline
Precondiciones & El medio de comunicación tiene que ser uno de los incluidos en la lista. Hay que tener conexión a internet para acceder a los RSS de estos medios. \\
\hline
Acciones & En la pestaña de lectura de RSS, el usuario elegirá en un desplegable el medio de comunicación del que quiere leer la noticia.

Pulsando el botón de ``Guardar Nuevas'', el usuario podrá almacenar las que no estuvieran ya en la base de datos. \\
\hline
Postcondiciones & Se mostrará una tabla con las últimas noticias del medio de comunicación seleccionado. \\
\hline
Excepciones & En caso de que no haya conexión a Internet, la lectura de noticias no se podrá realizar. \\
\hline 
Importancia & Alta \\
\hline 
\end{tabular}
\captionof{table}{Caso de uso 5}
\end{center}

\clearpage
\subsubsection{CU6 - Lectura de XML}

\begin{center}
\begin{tabular}{ | m{3cm} | m{10cm}| } 
\hline
Caso de uso & Lectura de XML \\ 
\hline
Versión & 1.0 \\ 
\hline
Autor & Miguel Rodríguez \\ 
\hline
Requisitos & RF - 6: Los usuarios podrán añadir un corpus de noticias entero, indicando el tema de debate y la etiqueta de cada noticia, a través de un archivo XML. \\ 
\hline
Descripción & El usuario ordena leer los documentos XML que contienen los corpus situados en el directorio especificado.\\
\hline
Precondiciones & Los documentos XML tienen que tener el formato determinado. \\
\hline
Acciones & En la pestaña de lectura de nuevos conjuntos, el usuario pulsará el botón de ``Leer''. 

Se cargará una nueva pestaña donde si el usuario pulsa ``Guardar Nuevas'', se almacenarán las noticias, con las etiquetas ya aplicadas de antemano. \\
\hline
Postcondiciones & Se guardarán las nuevas noticias con etiquetas ya aplicadas, al contrario que en los otros formatos. \\
\hline
Excepciones & En caso de que no hubiera archivos XML o que no tengan el formato correcto, la lectura no se realizará correctamente. \\
\hline 
Importancia & Alta \\
\hline 
\end{tabular}
\captionof{table}{Caso de uso 6}
\end{center}

\clearpage
\subsubsection{CU7 - Creación de Dataset}

\begin{center}
\begin{tabular}{ | m{3cm} | m{10cm}| } 
\hline
Caso de uso & Creación de dataset \\ 
\hline
Versión & 1.0 \\ 
\hline
Autor & Miguel Rodríguez \\ 
\hline
Requisitos & RF - 7: Los usuarios podrán crear nuevos temas de debate, escribiendo el nombre del propio tema y el nombre de dos clases, que se usarán posteriormente de etiquetas para ese tema. \\ 
\hline
Descripción & El usuario tiene la opción de introducir nuevos dataset para las noticias indicando su nombre y el de sus 2 posibles clases. \\
\hline
Precondiciones & Los temas de debate deben de tener 2 clases posibles. \\
\hline
Acciones & En la pestaña inicial, el usuario debe pulsar el botón de ``Añadir Nuevo Dataset''.

En la pestaña a la que ha sido redirigido, escribirá el nombre del nuevo dataset y el de sus dos clases, que serán las etiquetas que luego se podrán añadir a las noticias.

Al pulsar el botón de guardar cambios, se actualizará la base de datos con esta información. \\
\hline
Postcondiciones & Se añadirá a la lista de datasets el escrito en la aplicación, y a todas las noticias almacenadas también se les añadirán campos para soportar las nuevas etiquetas.\\
\hline
Excepciones & En caso de que no se rellenen los tres campos, no se guardará correctamente el nuevo dataset. \\
\hline 
Importancia & Alta \\
\hline 
\end{tabular}
\captionof{table}{Caso de uso 7}
\end{center}
\apendice{Especificación de diseño}

\section{Introducción}

En este apartado se analizará el diseño de la aplicación desde las diferentes perspectivas que componen el sistema. Se hablará del diseño de los datos, del flujo de ejecución de la aplicación y del diseño de la estructura de la misma.

\section{Diseño de datos}

Como ya se ha mencionado en la memoria, la base de datos que se usa en la aplicación es NoSQL, lo cual quiere decir que no tendrá un sistema relacional de tablas como es costumbre. 

También hay que destacar que, como la herramienta no está pensada para el uso abierto a cualquier usuario, si no a uno en concreto, no se consideró necesario hacer un sistema de registro de usuarios, con lo cual las colecciones/tablas habituales con la información de los usuarios aquí son inexistentes.

Sin embargo, si que existen dos colecciones (Se entiende por colección el equivalente a una tabla en SQL) en torno a las cuáles se estructura el almacenamiento de la información. A continuación describimos qué contienen cada una de ellas:

\begin{itemize}

\item \textbf{Noticias:} Es la colección que almacena los datos importantes de cada una de las noticias guardadas en la hemeroteca.

Los campos de esta colección son los siguientes:

\begin{itemize}

\item \textbf{Title}: El titular de la noticia.
\item \textbf{Author}: El autor de la noticia.
\item \textbf{Text}: El texto de la noticia.
\item \textbf{PublishDate}: La fecha de publicación de la noticia.
\item \textbf{Source}: La fuente de la noticia.
\item \textbf{Link}: Ubicación o url de la noticia.
\item \textbf{Tag}: El conjunto de etiquetas para cada dataset de la noticia.

\end{itemize}

\item \textbf{Datasets:} Es la colección que almacena los diferentes temas de debate sobre los que luego se podrá etiquetar y realizar predicciones con las noticias, siendo las etiquetas el nombre de las clases. 

Los campos que contiene son estos:

\begin{itemize}

\item \textbf{Dataset}: El nombre del tema de debate.
\item \textbf{Clase1}: El nombre de la primera clase.
\item \textbf{Title}: El nombre de la segunda clase.

\end{itemize}

\end{itemize}

\section{Diseño procedimental}

\section{Diseño arquitectónico}



\apendice{Documentación técnica de programación}

\section{Introducción}

En esta sección se explicarán todos los apartados necesarios para que, en el caso de que un programador externo entrara a trabajar en esta aplicación, pudiera familiarizarse rápido con ella y su desarrollo.

\section{Estructura de directorios}

\emph{Prototipo Flask} es el directorio en el que se ubican todos los archivos de la herramienta. En esta carpeta podemos encontrar los siguientes directorios:

\begin{itemize}

\item \textbf{app}: Es el directorio más importante de la propia aplicación, ya que contiene todos los archivos de \emph{back-end} (scripts de Python con la funcionalidad) y \emph{front-end} (documentos HTML para las pestañas) que forman la herramienta, junto con las carpetas con los estilos CSS, los fuentes de \emph{Bootstrap}, \emph{JQuery} y alguna biblioteca adicional, y las imágenes que se usan en las pestañas.

\item \textbf{datasets}: En este directorio es donde se deben de dejar los archivos XML que contengan corpus de noticias enteros. Este formato está pensado para casos en los que el usuario ya tuviera un conjunto de noticias con las clases asignadas de forma previa al uso de la herramienta.

\item \textbf{pdf}: En esta ubicación se deben de depositar los archivos PDF que sigan el formato indicado para su posterior lectura por parte de la aplicación.

\item \textbf{tests}: Aquí se encuentran los ficheros de código Python con las pruebas en Selenium para los casos de uso. Con la perspectiva de probarles por separado, se creó un archivo con una prueba para cada caso de uso.

\end{itemize}

Aparte de lo anterior, en el directorio raíz tenemos los siguientes archivos:

\begin{itemize}

\item \emph{config.py}: Fichero que contiene las variables de configuración (En este caso, las rutas a las carpetas donde se ubican los PDF, XML... etc)s.

\item \emph{requirements.txt}: Este fichero de texto plano contiene una lista de todos los recursos usados, para que cuando se instale la aplicación en otro equipo se puedan instalar a la vez usando este archivo en vez de uno a uno.

\item \emph{run.py}: Este código de Python es el que arranca la aplicación, haciendo uso del directorio \emph{app}.

\item \emph{setup.sh}: Contiene un conjunto de comandos simples, como el arranque del entorno virtual y de la base de datos, para que sólo sea necesario ejecutar este archivo antes de usar la aplicación.

\end{itemize}

\section{Manual del programador}

En este apartado se mencionarán los distintos puntos a tener en cuenta en relación al desarrollo de la aplicación.

\subsection{Minería de datos}

Todo lo relacionado con el aprendizaje automático de los textos de las noticias se encuentra en el script \emph{forms.py}, que se encuentra dentro de la carpeta app. 

En la clase \emph{Query}, que es la que se llama desde la aplicación para estas operaciones, tenemos dos métodos principales: \emph{classification} y \emph{explanation}. Si se quisiera modificar la predicción de noticias que se realiza en la tabla de resultados, habría que cambiar el código de la primera función, mientras que si se quieren hacer cambios en las explicaciones se deberá modificar la segunda función.

En lo que respecta al modelo, se refactorizó código extrayendo en una función apartada, llamada \emph{prepare model}, el preparado de los datos de entrenamiento a partir de las noticias.

A continuación se muestra una captura del código donde se puede observar esta estructura.

\imagen{MachineLearningCode}{Ubicación del código utilizado en el aprendizaje automático}

\subsection{Lectura de noticias}

Toda la funcionalidad relacionada con el almacenamiento y obtención de noticias, e interacción con la web se encuentra en el fichero \emph{views.py}. Allí, se puede encontrar una función para cada pestaña de la aplicación, ocupada de la funcionalidad \emph{back-end}. 

En la siguiente imagen se muestran estas funciones.

\imagen{NewsCode}{Funciones encargadas de la interacción con las pestañas}

\section{Compilación, instalación y ejecución del proyecto}

En los siguientes puntos se van a detallar qué elementos son imprescindibles para la configuración e instalación de la aplicación. Este proceso es muy similar al que se realiza en el Mega-Tutorial de Flask de Miguel Gringberg \cite{flaskmegatutorial}.

\subsection{Python}
Es esencial que en el equipo que vayamos a usar la herramienta esté Python instalado. Más concretamente, la versión 2.7 es en la que se ha desarrollado nuestro código.

En las últimas versiones de los Sistemas Operativos Ubuntu, ya viene una versión de Python 2 y Python 3 instalada, por lo tanto no requiere acciones extra por nuestra parte.

\subsection{Pip}

Otra herramienta muy útil a la hora de instalar bibliotecas de Python es Pip. Se puede instalar simplemente con el siguiente comando: 

\begin{minted}{bash}
sudo apt-get install python-pip
\end{minted}

\subsection{NLTK}

Para poder usar las palabras vacías, necesitamos descargarnos el corpus \emph{stopwords} de nltk, para ello, en un script Python, basta con los siguientes comandos:

\begin{minted}{bash}
import nltk
nltk.download()
\end{minted}

Se abrirá el menú de instalación de nltk, y desde ahí se podrá elegir que biblioteca queremos descargarnos (en este caso, la llamada ``stopwords'').

\subsection{Pymongo}

Antes de comenzar la instalación de la aplicación, hay que configurar en el equipo la base de datos NoSQL. Antes de la extensión de Python, necesitamos tener MongoDB en nuestro equipo, lo cual conseguiremos siguiendo los pasos que indican en la guía de la página oficial de MongoDB \cite{Chodorow:2010:MDG:1941134}, en este enlace: \url{https://docs.mongodb.com/v3.2/tutorial/install-mongodb-on-ubuntu/}.

Para instalar Pymongo, teniendo Pip en el equipo basta con escribir lo siguiente:

\begin{minted}{bash}
sudo python -m pip install pymongo
\end{minted}


Para activar el servicio de MongoDB en el Sistema Operativo sólo hace falta escribir en la terminal:

\begin{minted}{bash}
sudo service mongod start
\end{minted}

\subsection{Entorno Virtual}

Para poder usar Flask, se creará un entorno virtual en el que se va a realizar la instalación de todo lo relacionado con Flask y se ejecutará la aplicación.

Para poder crear este entorno virtual, lo primero que hay que hacer es instalar \emph{virtualenv}, con el siguiente comando:

\begin{minted}{bash}
sudo apt-get install python-virtualenv
\end{minted}

\subsection{Flask}

Una vez esté instalado \emph{virtualenv}, desde la ventana de comandos habrá que situarse en la ubicación del directorio raíz del proyecto y ejecutar el siguiente comando:

\begin{minted}{bash}
virtualenv flask
\end{minted}


Aprovechando que hay un archivo llamado \emph{requirements} que contiene todas las bibliotecas que se han usado en el desarrollo, basta con el siguiente comando para que se instalen todas en orden sin necesidad de instalarlas una a una:

\begin{minted}{bash}
sudo flask/bin/pip install -r requirements.txt
\end{minted}


Una vez está instalado todo lo necesario, hay que modificar los permisos del fichero \emph{run.py} para que nos permita ejecutarle. Se hace de la siguiente manera:

\begin{minted}{bash}
sudo chmod a+x run.py
\end{minted}

Una vez sea accesible este archivo, ya se podría ejecutar la aplicación. Es crucial mencionar que la primera vez que se ejecuta la aplicación en un ordenador nuevo, no va a tener ningún dataset, y por tanto se debe usar el script \emph{createdataset.py}, ubicado en el directorio raíz, para crear los dataset iniciales con los que poder empezar a etiquetar noticias.


\section{Pruebas del sistema}

\subsection{Pruebas de funcionalidad}

Aparte de las pruebas convencionales de funcionamiento realizadas por usuarios, se han creado un conjunto de pruebas simples automatizadas con la herramienta \emph{Selenium}. 

No obstante, al contrario de su uso habitual, que consistiría en grabar un conjunto de pasos con el plugin del navegador y exportarlo en algún lenguaje como Java, hemos usado las bibliotecas de \emph{Selenium} para Python, evitando así usar diferentes lenguajes de programación, aún a costa de no usar el plugin para Mozilla Firefox.

En la carpeta \emph{tests} del proyecto, se pueden encontrar un conjunto de scripts numerados del 1 al 7, que contienen una prueba simple para cada caso de uso. Para evitar hacer cambios innecesarios en la base de datos, se ha evitado ejecutar la funcionalidad de guardar nuevas noticias en las pruebas, ya que si por ejemplo, en el caso de uso de lectura de RSS, se estuvieran haciendo pruebas constantes, podría llenarse la base de datos de noticias que no se querían almacenar en un primer momento.

Además de esto, hay un script auxiliar llamado \emph{deletedataset} que se puede usar para borrar un conjunto de datos de la aplicación. El motivo de que esté guardado es que, realizando pruebas, se acababan almacenando demasiados conjuntos de pruebas innecesarios en la aplicación, y esto podía afectar al rendimiento. Por tanto, ejecutando este script con el nombre del conjunto que queramos borrar ahorra bastante trabajo al programador que esté probando la herramienta.

\subsection{Calidad del código}
Además de las pruebas de funcionalidad, se llevó a cabo una evolución progresiva de la calidad del código usando la herramienta \emph{Pylint}, que detecta todos los errores, avisos y mejoras de refactorización que se pueden realizar en programas escritos en Python. Esta herramienta se usó tanto en \emph{views.py} como en \emph{forms.py}, que son los dos scripts que reunen la mayoría de la funcionalidad de la aplicación.

Por ejemplo, en las siguientes imágenes podemos ver la evolución de calidad del código del fichero \emph{views} en función de la nota recibida por \emph{Pylint}:

\imagen{PreviousQuality}{Nota obtenida antes de realizar refactorizaciones}

\imagen{NewQuality}{Nota obtenida después de varias semanas, aplicando algunas técnicas de refactorización}

\imagen{NewQuality2Views}{Nota obtenida en los últimos días de desarrollo}

\imagen{NewMessages}{Algunos de los factores que valora Pylint}

Del mismo modo, podemos observar de forma breve la evolución de \emph{forms} a lo largo de las semanas, aplicando algunas técnicas de refactorización:

\imagen{FormsQuality1}{Nota obtenida antes de aplicar cambios}

\imagen{FormsQuality3}{Nota después de los primeros cambios}

\imagen{FormsQuality4}{Nota después de haber aplicado la mayoría de los cambios}


Cabe destacar que, aunque sigue habiendo margen de mejora, era imposible obtener un 10/10 con esta herramienta, pues es conflictiva en algunos aspectos. Por poner un ejemplo, evalúa con mala nota las líneas con más de 100 columnas de código, pero si se introduce un salto de línea para dividir la longitud, penaliza que haya un salto de línea inapropiado.
\apendice{Documentación de usuario}

\section{Introducción}
En este apartado se explicarán las condiciones que tiene que cumplir el equipo en el que se ejecute la aplicación, además de mostrar visualmente cómo se debe usar la misma.

\section{Requisitos de usuarios}

\subsection{Software}

Aunque en el apartado ``Manual del programador'' se explica con más detalle qué programas y librerías son necesarias para ejecutar la aplicación, los resumiremos en un listado breve como el siguiente:

\begin{itemize}

\item Sistema Operativo: La herramienta se ha desarrollado y está pensada para utilizarse en Ubuntu y, más específicamente, la versión 16.

\item Navegador Web: Durante el desarrollo se utilizó Mozilla Firefox, que además es el que usa \emph{Selenium} para realizar las pruebas. No obstante, la herramienta ha sido probada en otros navegadores como el incluido en Ubuntu por defecto, y funciona correctamente.

\item Python: Ya viene incluido en Ubuntu. Si aun así se deseara descargar, habría que elegir la versión 2.7.

\item Pip: Para realizar la instalación de la mayoría de bibliotecas que usa la aplicación.

\item MongoDB: Necesario para poder montar una base de datos con las noticias.

\item Virtualenv: El entorno virtual en el que se instalarán todo y desde donde se ejecutará la aplicación.

\item NLTK: Al menos, el corpus correspondiente a las palabras vacías.
\end{itemize}

\subsection{Hardware}

En lo relacionado con los aspectos técnicos del equipo en el que se ejecute la aplicación, no hay que tener en cuenta requisitos especialmente exigentes. Si que hay que contemplar que, a mejor \emph{hardware} esté disponible, las operaciones más costosas se realizarán con más facilidad.

Por tomar un punto de referencia mínimo, vamos a citar los requisitos hardware de Mozilla Firefox para los Sistemas Operativos convencionales \cite{mozillarequirements}:

\begin{itemize}

\item Procesador: Pentium 4 o superior

\item Tamaño de Memoria RAM: 512MB (Si se va a ejecutar en una máquina virtual, habrá que controlar que el equipo tenga suficiente RAM asignada a la misma para que funcione correctamente).

\end{itemize}

En términos de almacenamiento, se recomienda dejar libre alrededor de 1 GB para poder instalar varias de las bibliotecas adicionales que vienen incluidas en \emph{requirements}. No ocupan tanto espacio, pero lo más aconsejable es dejar un margen por si se implementaran nuevas funcionalidades o se llegaran a instalar nnuevas bibliotecas de terceros.

\section{Instalación}

El conjunto de elementos que hay que instalar para poder ejecutar la aplicación (Python, MongoDB ... etc) viene explicado con detalle en el apartado ``Manual del Programador'', pues los pasos a seguir son los mismos.

Habiendo realizado todas las acciones mencionadas antes, lo único necesario ahora sería ejecutar el archivo \emph{setup.sh} que hay en el directorio raíz, con el siguiente comando:

\begin{minted}{bash}
sudo bash setup.sh
\end{minted}

Este archivo simplemente reúne el conjunto de comandos que harían falta para ejecutar la herramienta, que consiste en:

\begin{itemize}

\item Arrancar base de datos.

\item Preparar entorno virtual.

\item Ejecutar \emph{run.py}

\end{itemize} 

De esta manera, solo hay que ejecutar un comando en vez de varios. Una vez escrito, la aplicación estará preparada y entrando en la url que se indica por la terminal, se puede acceder a la aplicación.

\section{Manual del usuario}

En este manual se van a explicar cómo se realizan las diferentes acciones que se pueden llevar a cabo dentro de la aplicación.

\subsection{Consulta de noticias}

Estando en la página de inicio, como se puede ver en la imagen \ref{fig:CapturaIndex}, hay que introducir la palabra clave que queramos buscar en el cuadro de texto, y pulsar en el botón ``Buscar''.

\imagen{CapturaIndex}{Ejemplo de palabra clave de búsqueda en la herramienta}

Tras un tiempo de espera, aparecerá en pantalla una tabla como la de la imagen \ref{fig:CapturaIndex2}

\imagen{CapturaIndex2}{Tabla de resultados}

\subsection{Etiquetado y Explicación}

Una vez tengamos acceso a la tabla de resultados de la búsqueda, podemos realizar dos acciones con las noticias, siempre respecto al dataset seleccionado en el desplegable que hay encima de la tabla:

\begin{itemize}

\item Primero, podemos realizar un etiquetado manual de las noticias. Para ello, basta con pulsar el botón "Etiquetar" de la noticia para que en la celda de la etiqueta aparezcan opciones de selección única, como se ve en la imagen \ref{fig:CapturaIndex2}. Cuando hayamos etiquetado las noticias que queramos, pulsamos en ``Guardar Resultados'' para aplicar los cambios.

\item La otra opción sería ver una explicación de la predicción de la noticia, pulsando en el botón ``Clasificar'', que abrirá otra pestaña en la que se puede ver una explicación de los criterios para elegir una clase como en la imagen \ref{fig:CapturaExplicacion}.

\imagen{CapturaExplicacion}{Explicación para la predicción de la clase de una noticia}
\end{itemize}

\subsection{Lectura de Noticias}

Si, por el contrario, lo que se quiere es aumentar el corpus de noticias en la base de datos, hay varias opciones de hacerlo.

\subsubsection{Lectura de PDF}

Pulsando en el apartado ``Lectura PDF'' de la barra de navegación, accedemos a la pestaña de lectura de documentos PDF. Allí, se indica al usuario en qué directorio debe dejar los archivos. La interfaz es la que se puede ver en \ref{fig:PDFpage}.

\imagen{PDFpage}{Pestaña de lectura de PDF}

Pulsando en el botón, nos llevará a una tabla de resultados igual que la de \ref{fig:PDFread}.

\imagen{PDFread}{Resultados de la lectura de PDF}

Y si allí se pulsa en ``Guardar Nuevas'', se almacenarán las noticias que no estuvieran ya en la base de datos.

Es crucial mencionar que las noticias se deben guardar en los PDF en un formato concreto, de lo contrario no se podrá realizar la lectura de la noticia correctamente. El formato es:

\begin{itemize}

\item \textbf{(Titular) *titular de la noticia* (Titular)}  

\item \textbf{(Autor) *autor de la noticia* (Autor)}

\item \textbf{(Fecha) *fecha de la noticia* (Fecha)}

\item \textbf{(Cuerpo) *texto de la noticia* (Cuerpo)}

\end{itemize}

Esta decisión se tomó debido a que inicialmente había numerosos PDF con formatos distintos de los que leer noticias, y no se podía programar una solución satisfactoria para todos los casos. En la carpeta \emph{pdf} de la aplicación, hay un documento de ejemplo con el formato adecuado.

\subsubsection{Lectura de RSS}

Pulsando en el apartado ``RSS Scraping'' de la barra de navegación, accedemos a una pestaña que muestra un desplegable con los nombres de varios medios de comunicación, como se puede ver en \ref{fig:RSS1}. Si elegimos uno de ellos y hacemos click en ``Aceptar'', nos mostrará una tabla con las últimas noticias en el RSS de ese medio de comunicación.

\imagen{RSS1}{Imagen inicial de la pestaña de lectura de RSS}

Tras un breve tiempo de cargar, aparecerán los resultados en una tabla como la de \ref{fig:RSS2}.

\imagen{RSS2}{Resultados de lectura de RSS}

Igual que con el resto de formatos, pulsando en ``Guardar Nuevas'' almacenaremos en la base de datos las que no estén ya guardadas.

\subsubsection{Lectura de Dataset}

Este apartado está pensado para introducir corpus de noticias enteros que el usuario pudiera tener ya preparados antes de haber usado la aplicación, por medio de un archivo XML. En esta pestaña, de forma similar a la de lectura de PDF, se indica cuál es la carpeta en la que hay que dejar los archivos, como se puede observar en \ref{fig:Corpus1}.

\imagen{Corpus1}{Pestaña de lectura de Corpus}

El formato que deben de tener los XML es el siguiente:

\imagen{Corpus}{Estructura de los corpus en XML}

Cuando se hayan leído las noticias del corpus, se devolverá en una tabla de resultados similar al resto, pero en esta ocasión las noticias ya tendrán una etiqueta asignada. Esta tabla es como la de la imagen \ref{fig:Corpusread}

\imagen{Corpusread}{Tabla de resultados con el corpus}

Al igual que en los otros apartados, se pueden guardar las nuevas noticias pulsando el botón correspondiente.

\subsection{Creación de nuevos Conjuntos de datos}

Desde la pestaña de inicio, si pulsamos el botón ``Crear nuevo Dataset'', nos llevará a una nueva ventana donde podemos introducir en los cuadros de texto el nombre del conjunto y de las 2 clases que lo van a formar, y añadirle a la lista de los existentes.

\imagen{NewDataset}{Pestaña de creación de Datasets}

Una vez creado, en la tabla de resultados de noticias podrán realizarse etiquetados para ese conjunto de datos.

\subsection{Estadísticas de las noticias}

La herramienta permite ver, de forma básica, un conjunto de gráficos para controlar la cantidad de noticias que hay almacenadas en relación a varios factores. En la pestaña de inicio, si se pulsa ``Ver estadísticas'', aparecerá una vista en la que podemos elegir si ver la cantidad de noticias en función del autor, de la fecha, o de la fuente, y mostrará el gráfico correspondiente.

No es una funcionalidad con mucha profundidad pero el usuario puede considerar interesante saber este tipo de información para añadir datos cuando realice estudios sobre las noticias. Los gráficos tienen el aspecto de la imagen \ref{fig:Statistics}

\imagen{Statistics}{Gráfico de noticias por autor}

\subsection{Consideraciones adicionales}

Al usar la herramienta en un nuevo ordenador, hay que tener en cuenta que no habrá ni noticias ni conjuntos de datos almacenados en la base de datos. 

Se podría solucionar este problema usando las opciones que da la propia herramienta, pero para facilitar la toma de contacto, hay un script llamado \emph{createdatasets} que añade a la base de datos los dos conjuntos iniciales con los que se ha trabajado, y en la carpeta \emph{dataset} está el corpus facilitado por nuestra colaboradora María Isabel Menéndez. La idea es que con estos dos sencillos pasos ya haya una base sobre la que seguir trabajando con las funcionalidades de la aplicación.

\subsection{Máquina virtual del programador}

Si, por el contrario, no se desea seguir el proceso de instalación indicado en el manual del programador, y poder pasar directamente a la ejecución de la aplicación, se puede tomar la alternativa de descargar la máquina virtual en la que se ha realizado el desarrollo de la aplicación, en la que está todo lo necesario ya instalado.

El enlace para descargar esta máquina virtual es el siguiente \url{https://mega.nz/#!RQEGkLLb!ybxP6ZS2HmnhrWxEqfJtHQi8ViBVeaRCvo8iStOnMHg}. Una vez descargada, sólo será necesario importarla en \emph{VirtualBox} y arrancarla. Para ejecutar la aplicación, en la terminal, dentro de la carpeta de proyecto, habrá que escribir este comando:

\begin{minted}{bash}
sudo bash setup.sh
\end{minted}

El usuario de la máquina virtual se llama \emph{alumnotfg} y la contraseña es \textbf{alumnotfg2017}.


\bibliographystyle{plain}
\bibliography{bibliografiaAnexos}

\end{document}
